\documentclass[a4paper,12pt]{article}
\usepackage[utf8]{inputenc}
\usepackage[spanish]{babel}
\usepackage{amsmath}
\usepackage{amssymb}
\usepackage{amsthm}
\usepackage{geometry}
\usepackage{fancyhdr}
\usepackage{tcolorbox}
\usepackage{booktabs}
\usepackage{multicol}
\usepackage{float}

\geometry{top=2.5cm, bottom=2.5cm, left=3cm, right=3cm}

% Configuración de encabezado
\setlength{\headheight}{14.49998pt}
\pagestyle{fancy}
\fancyhf{}
\lhead{Licenciatura en Sistemas de la Información}
\rhead{Unidad 1: Sucesiones Y Series}
\cfoot{\thepage}

\title{\textbf{Resumen Unidad 1: Sucesiones Numéricas y Series de Funciones}}
\author{Ecuaciones Diferenciales y Cálculo Multivariado}
\date{}

\begin{document}

\maketitle

\part{Sucesiones}

\section{Sucesiones Numéricas}

\subsection{Definición}
Una sucesión numérica se define de dos formas equivalentes:
\begin{enumerate}
    \item Como un \textbf{conjunto ordenado de números} afectados por un índice natural que indica su posición: $\{a_1, a_2, a_3, ..., a_n, ...\}$.
    \item Como una \textbf{función} cuyo dominio son los números naturales ($\mathbb{N}$) y su rango los números reales ($\mathbb{R}$):
    $$ f: \mathbb{N} \rightarrow \mathbb{R} \quad / \quad \forall n \in \mathbb{N}: a_n = f(n) $$
\end{enumerate}

\subsection{Formas de Representación}
\begin{itemize}
    \item \textbf{Término n-ésimo ($a_n$):} Fórmula explícita. Ej: $a_n = \frac{n}{n^2+1}$.
    \item \textbf{Coloquial:} Descripción verbal. Ej: "1 si es par, 0 si es impar".
    \item \textbf{Regla de Recurrencia:} El término depende de los anteriores. 
    Ej: $a_n = 3a_{n-1} - a_{n-2}$, dados $a_1$ y $a_2$.
    \item \textbf{Geométrica:} Puntos en el plano cartesiano $(n, a_n)$ o sobre la recta real.
\end{itemize}

\section{Sucesiones Especiales}

\subsection{Sucesión Aritmética}
Cada término se obtiene sumando al anterior una constante $s$ (diferencia).
$$ a_n = a + (n-1)s $$
\textit{Ejemplo:} $a_n = 2 + (n-1)\frac{1}{3}$.

\subsection{Sucesión Geométrica}
Cada término se obtiene multiplicando al anterior por una constante $q$ (razón).
$$ a_n = a \cdot q^{n-1} $$
\textit{Ejemplo:} $a_n = 2\left(\frac{1}{2}\right)^{n-1}$.

\subsection{Otras}
\begin{itemize}
    \item \textbf{Aritmética-Geométrica:} $a_n = [a + (n-1)s] \cdot q^{n-1}$.
    \item \textbf{Armónica o recíproca:} $a_n = \frac{1}{n}$.
\end{itemize}

\subsection{Operaciones con sucesiones}
Dadas dos sucesiones numéricas $\{a_n\}$ y $\{b_n\}$, $\alpha, \beta$ dos constantes aribitrarias:
\begin{itemize}
    \item El \textbf{producto} $\{a_n \cdot b_n\}$ es también una sucesión. ${(ab)_n}$
    \item La \textbf{combinación lineal} de ${\alpha a_n + \beta b_n}$ es también una sucesión. ${(\alpha a_n + \beta b_n)}$
    \item ${b_n} \neq 0, \forall n \in \mathbb{N} \Rightarrow \left\{\frac{1}{b_n}\right\}$ es también una sucesión. $\left\{(\frac{1}{b})_n\right\}$
    \item Si ${b_n} \neq 0, \forall n \in \mathbb{N} \Rightarrow \left\{\frac{a_n}{b_n}\right\}$ es también una sucesión. $\left\{(\frac{a}{b})_n\right\}$
\end{itemize}

\begin{tcolorbox}[title=Igualdad de Sucesiones,colback=blue!5!white,colframe=blue!75!black]
    $${a_n} = {b_n} \Leftrightarrow a_j = b_j, \forall j \in \mathbb{N}$$
\end{tcolorbox}

\section{Monotonía y Acotación}

\subsection{Sucesiones Monótonas}
Una sucesión es \textbf{monótona} si cumple con alguna de las siguientes condiciones:
\begin{itemize}
    \item \textbf{Creciente:} $\forall n, a_n < a_{n+1}$.
    \item \textbf{Decreciente:} $\forall n, a_n > a_{n+1}$.
    \item \textbf{No decreciente:} $a_n \le a_{n+1}$.
    \item \textbf{No creciente:} $a_n \ge a_{n+1}$.
\end{itemize}

\subsection{Sucesiones Acotadas}
Una sucesión $\{a_n\}$ está \textbf{acotada por arriba} si existe un $M \in \mathbb{R}$ tal que:
$$ \forall n \in \mathbb{N}, a_n \leqslant M $$
y \textbf{acotada por abajo} si existe un $m \in \mathbb{R}$ tal que:
$$ \forall n \in \mathbb{N}, a_n \geqslant m $$
Si una sucesión está acotada por arriba y por abajo, se dice que está \textbf{acotada}.

\subsection{Teorema de Bolzano}
Toda sucesión monótona acotada tiene un límite finito.

\section{Límite y Convergencia}

\subsection{Definición de Límite}
Una sucesión tiene límite $L$ (es convergente) si:
$$ \lim_{n \to \infty} a_n = L \iff \forall \epsilon > 0, \exists N > 0 \ / \ \forall n > N \implies |a_n - L| < \epsilon $$

\subsection{Clasificación según el límite}
\begin{itemize}
    \item \textbf{Convergente:} Tiene límite finito único.
    \item \textbf{Divergente:} No tiene límite finito (puede tender a $\infty$ u oscilar indefinidamente).
\end{itemize}

\subsection{Teoremas de Convergencia}
\begin{enumerate}
    \item Si $\lim_{n \to \infty} a_n = L$ y la función $f$ es continua en $L$, entonces: 
    $$ \lim_{n \to \infty} f(a_n) = f(L) $$
    \item La sucesión ${r^n}$ \textbf{converge} si $-1 < r \leqslant 1$ y \textbf{diverge} para el resto de valores.
    $$\lim_{n \to \infty} r^n = \begin{cases} 0 & \text{si } -1 < r \leqslant 1 \\ 1 & \text{si } r = 1 \end{cases}$$
\end{enumerate}

\subsection{Teoremas de Convergencia: Monotonía y Acotación}
\begin{enumerate}
    \item \textbf{Monotonía y Acotación:}
    \begin{itemize}
        \item Si es monótona creciente y tiene cota superior $\rightarrow$ \textbf{Converge}.
        \item Si es monótona decreciente y tiene cota inferior $\rightarrow$ \textbf{Converge}.
    \end{itemize}
    \textbf{\textit{Toda sucesión monótona y acotada es convergente.}}

    \item \textbf{Condición Necesaria (pero no suficiente):}
    \begin{itemize}
        \item Si $\{a_n\}$ converge $\implies$ $\{a_n\}$ está acotada.
        \item \textbf{Contrarrecíproco:} Si no está acotada $\implies$ es \textbf{divergente}.
        \item \textit{Nota:} La acotación por sí sola no garantiza convergencia (ej: $(-1)^n$ es acotada pero oscila).
    \end{itemize}
\end{enumerate}

\subsection{Leyes de los límites para sucesiones}
\begin{enumerate}
    \item $\lim_{n \to \infty} (a_n + b_n) = \lim_{n \to \infty} a_n + \lim_{n \to \infty} b_n$
    \item $\lim_{n \to \infty} (a_n - b_n) = \lim_{n \to \infty} a_n - \lim_{n \to \infty} b_n$
    \item $\lim_{n \to \infty} (C \cdot a_n) = C \cdot \lim_{n \to \infty} a_n$, para cualquier constante $C$
    \item $\lim_{n \to \infty} (a_n \cdot b_n) = \lim_{n \to \infty} a_n \cdot \lim_{n \to \infty} b_n$
    \item $\lim_{n \to \infty} \left(\frac{a_n}{b_n}\right) = \frac{\lim_{n \to \infty} a_n}{\lim_{n \to \infty} b_n}$, si $\lim_{n \to \infty} b_n \neq 0$
    \item $\lim_{n \to \infty} (a_n)^r = \left(\lim_{n \to \infty} a_n\right)^r = \lim_{n \to \infty} r^{a_n} = r^{\lim_{n \to \infty} a_n}$, para $r \in \mathbb{R}$ y $r > 0$ y $a_n > 0$
\end{enumerate}

\begin{tcolorbox}[title=Teorema de compresión,colback=blue!5!white,colframe=blue!75!black]
    Si $a_n \leqslant b_n \leqslant c_n \forall n \geqslant n_0$ y $\lim_{n \to \infty} a_n = \lim_{n \to \infty} c_n = L \Rightarrow \lim_{n \to \infty} b_n = L$.
\end{tcolorbox}
%Agregar a mano lim del valor absoluto (p698)

\begin{tcolorbox}[title=Teorema de compresión,colback=blue!5!white,colframe=blue!75!black]
    Al hallar una indeterminación en el límite de una sucesión, no puede usarse L'Hôpital. Sin embargo, se puede aplicar a la función relacionada:
    $$ \lim_{n \to \infty} \frac{\ln n}{n} = \lim_{x \to \infty} \frac{\ln x}{x} = \lim_{x \to \infty} \frac{1/x}{1} = 0 $$
\end{tcolorbox}

\section{Infinitésimos y Sucesiones Infinitas}

\subsection{Sucesiones Infinitésimas}
Son aquellas cuyo límite es 0.
$$ \lim_{n \to \infty} a_n = 0 $$
\textbf{Propiedad fundamental:} Un número $L$ es límite de $a_n \Leftrightarrow a_n = L + \alpha_n$, donde $\{\alpha_n\}$ es una sucesión infinitésima y $\alpha_n$ su enésimo elemento.

\subsection{Propiedades de Infinitésimos}
\begin{itemize}
    \item Suma de infinitésimos = Infinitésimo.
    \item Producto de infinitésimo $\times$ Acotada = Infinitésimo.
\end{itemize}

\subsection{Sucesiones Infinitas}
Aquellas donde $\lim_{n \to \infty} a_n = \infty$. Si $a_n$ es infinita ($a_n \neq 0$), entonces $\{\frac{1}{a_n}\}$ es infinitésimo.

\section{Álgebra de Límites}
Si $\{a_n\} \to A$ y $\{b_n\} \to B$:
\begin{align*}
    \lim_{n \to \infty} (a_n \pm b_n) &= A \pm B \\
    \lim_{n \to \infty} (C \cdot a_n) &= C \cdot A \\
    \lim_{n \to \infty} (a_n \cdot b_n) &= A \cdot B \\
    \lim_{n \to \infty} \left(\frac{a_n}{b_n}\right) &= \frac{A}{B} \quad (si \ B \neq 0) \\
    \lim_{n \to \infty} (a_n)^r &= A^r
\end{align*}

\newpage

\part{Series Infinitas}

\section{Definiciones Fundamentales}

\subsection{Definición de Serie}
Una serie infinita es la suma de los términos de una sucesión infinita $\{a_n\}_{n=1}^{\infty}$ y se denota como:
$$\sum_{n=1}^{\infty} a_n = a_1 + a_2 + a_3 + \dots + a_n + \dots$$

\subsection{Sumas Parciales y Convergencia}
Para determinar si una serie tiene una suma finita, se analiza la \textbf{sucesión de sumas parciales} $\{s_n\}$:
\begin{itemize}
    \item $s_1 = a_1$
    \item $s_2 = a_1 + a_2$
    \item $s_n = \sum_{i=1}^{n} a_i$
\end{itemize}

\begin{tcolorbox}[title=Definición de Convergencia,colback=blue!5!white,colframe=blue!75!black]
Si la sucesión $\{s_n\}$ es convergente y $\lim_{n \to \infty} s_n = s$ (donde $s$ es un número real), la serie $\sum a_n$ es \textbf{convergente} y su suma es $s$. Si el límite no existe o es infinito, la serie es \textbf{divergente}.
\end{tcolorbox}

\section{Propiedades}
\begin{itemize}
    \item La convergencia o divergencia de una serie no es afectada:
    \begin{itemize}
        \item al cambiar un número finito de términos.
        \item al multiplicar todos los términos por una constante: $\sum_{n=1}^{\infty} a_n  C = C \sum_{n=1}^{\infty} a_n$
    \end{itemize} 
    \item Dos series convergentes pueden sumar o restarse término a término y la serie resultante también converge y su suma $s$ es obtenida sumando o restando las sumas de las series dadas:
    $$\sum_{n=1}^{\infty} a_n \pm \sum_{n=1}^{\infty} b_n = \sum_{n=1}^{\infty} (a_n \pm b_n)$$
    \item Si se suma una serie convergente y una divergente, la serie resultante es divergente.
\end{itemize}

\section{Series Especiales}

\subsection{Serie Geométrica}
Se define como $\sum_{n=1}^{\infty} ar^{n-1}$. 
\begin{itemize}
    \item \textbf{Converge} si $|r| < 1$, su suma es $s = \frac{a}{1-r}$.
    \item \textbf{Diverge} si $|r| \geqslant 1$.
\end{itemize}
\textit{Resumen:} la suma de una serie geométrica convergente es: $\frac{\text{primer término}}{1 - \text{razón común}}$

\subsection{Serie p}
Se define como $\sum_{n=1}^{\infty} \frac{1}{n^p}$.
\begin{itemize}
    \item \textbf{Converge} si $p > 1$.
    \item \textbf{Diverge} si $p \leqslant 1$.
\end{itemize}

\subsection{Serie Armónica}
Es el caso especial de la serie $p$ donde $p=1$: $\sum_{n=1}^{\infty} \frac{1}{n}$.
\textbf{Siempre es divergente.}

\section{Convergencia Absoluta y Condicional}
\begin{itemize}
    \item \textbf{Convergencia Absoluta:} $\sum a_n$ es absolutamente convergente si $\sum |a_n|$ es convergente. En este caso, $\sum a_n$ es convergente.
    \item \textbf{Convergencia Condicional:} Una serie es condicionalmente convergente si es convergente pero $\sum |a_n|$ diverge.
\end{itemize}

\section{Pruebas de Convergencia}

\subsection{Prueba de la Divergencia (Término n-ésimo)}
Si $\lim_{n \to \infty} a_n$ no existe o si $\lim_{n \to \infty} a_n \neq 0$, entonces la serie $\sum a_n$ es \textbf{divergente}.
\begin{tcolorbox}
    Si $\sum_{n=1}^{\infty} a_n$ es \textbf{convergente}, entonces $\lim_{n \to \infty} a_n = 0$, \textit{\textbf{pero el recíproco no es necesariamente cierto}}.
\end{tcolorbox}

\subsection{Prueba de la Integral}
\begin{tcolorbox}
    La serie $\sum a_n$ \textbf{ converge } $\Leftrightarrow$ \textbf{ la integral impropia } $\int_{1}^{\infty} f(x) dx$ converge.
\end{tcolorbox}
Se aplica si $a_n = f(n)$, donde $f$ es continua, positiva y decreciente en $[1, \infty)$.

\subsubsection{Estimación del residuo para la prueba de la ingtegral}
Sea $R_n = s - s_n$ el residuo (error) al usar la suma parcial $s_n$ para aproximar la suma de la serie. Entonces:
$$ \int_{n+1}^{\infty} f(x) dx \leqslant R_n \leqslant \int_{n}^{\infty} f(x) dx $$

\subsection{Pruebas de Comparación}
\begin{enumerate}
    \item \textbf{Comparación Directa:} 
        \begin{itemize}
            \item Si $0 \leqslant a_n \leqslant b_n$ y $\sum b_n$ converge $\implies \sum a_n$ converge.
            \item Si $a_n \geqslant b_n \geqslant 0$ y $\sum b_n$ diverge $\implies \sum a_n$ diverge.
        \end{itemize}
    \item \textbf{Comparación del Límite:} Si $\lim_{n \to \infty} \frac{a_n}{b_n} = c$:
        \begin{itemize}
            \item Si $0 < c < \infty$, entonces $\sum a_n$ y $\sum b_n$ \textbf{convergen o divergen juntas}.
            \item Si $c = 0$ y $\sum b_n$ converge $\implies \sum a_n$ converge.
            \item Si $c = \infty$ y $\sum b_n$ diverge $\implies \sum a_n$ diverge.
        \end{itemize}
\end{enumerate}

\subsubsection{Estimación del residuo para la prueba de comparación}
Sea $R_n = s - s_n$ el residuo de $\sum a_n$ y $T_n = t - t_n$ el residuo de $\sum b_n$. Si $a_n \leqslant b_n$ para $n \geqslant N$, entonces $R_n \leqslant T_n$ para $n \geqslant N$.

\subsection{Prueba de la Serie Alternante - Criterio de Leibniz}
Para series de la forma $\sum (-1)^{n-1} b_n$ con $b_n > 0$. Converge si:
\begin{itemize}
    \item $b_{n+1} \leqslant b_n$ para toda $n$.
    \item $\lim_{n \to \infty} b_n = 0$.
\end{itemize}

\begin{tcolorbox}
    Si en una serie alternante $a_n > a_{n+1}$ pero $\lim_{n \to \infty} a_n \neq 0$, entonces la serie es  oscilante, por lo que no converge.\newline
    Si $\lim_{n \to \infty} a_n = 0$ pero $a_n$ no es decreciente, la serie diverge.
\end{tcolorbox}

\subsubsection{Estimación del residuo para la prueba de la serie alternante}
Sea $s$ la suma de la serie alternante y $s_n$ la suma parcial de los primeros $n$ términos. Entonces el residuo $R_n = s - s_n$ satisface:
$$ |R_n| \leqslant b_{n+1} $$

\subsection{Prueba de la Razón (D'Alembert)}
Sea $L = \lim_{n \to \infty} | \frac{a_{n+1}}{a_n} |$:
\begin{itemize}
    \item Si $L < 1$, es \textbf{absolutamente convergente}.
    \item Si $L > 1$, es \textbf{divergente}.
    \item Si $L = 1$, la prueba no concluye.
\end{itemize}

\subsection{Prueba de la Raíz (Cauchy)}
Sea $L = \lim_{n \to \infty} \sqrt[n]{|a_n|}$:
\begin{itemize}
    \item Si $L < 1$, es \textbf{absolutamente convergente}.
    \item Si $L > 1$, es \textbf{divergente}.
    \item Si $L = 1$, la prueba no concluye.
\end{itemize}

\subsection{Criterio de Raabe}
Sea $a_n > 0$ y $L = \lim_{n \to \infty} n \left( \frac{a_n}{a_{n+1}} - 1 \right)$:
\begin{itemize}
    \item Si $L > 1$, la serie $\sum a_n$ es \textbf{absolutamente convergente}.
    \item Si $L < 1$, la serie $\sum a_n$ es \textbf{divergente}.
    \item Si $L = 1$, la prueba no concluye.
\end{itemize}

%%---------------------------------------------------------------------------------------------

\section{Estrategia para probar series}
\begin{enumerate}
    \item \textbf{Serie p }($\sum \frac{1}{n^p}$): converge si $p > 1$ y diverge si $p \leqslant 1$.
    \item \textbf{Serie Geométrica }($\sum ar^{n-1}$ o $\sum ar^n$): converge si $|r| < 1$ y diverge si $|r| \geqslant 1$.
    \item Si la serie tiene forma similar a una serie p o geométrica, usar la \textbf{Prueba de Comparación} o la \textbf{Prueba del Límite de Comparación}. Si la serie tiene términos negativos, evaluar $\sum |a_n|$.
    \item Si es fácil ver que $\lim_{n \to \infty} a_n \neq 0$, aplicar la \textbf{Prueba de la Divergencia}.
    \item Si la serie es de forma $\sum (-1)^{n-1} b_n$ o $\sum (-1)^n b_n$, usar la \textbf{Prueba de la Serie Alternante}.
    \item Para las series con factoriales u otros productos, utilizar la \textbf{Prueba de la Razón}. Es útil pensar que $|\frac{a_{n+1}}{a_n}| \to 1$ cuando $n \to \infty$ para todas las series p y funciones racionales o algebraicas de $n$.
    \item Para la forma $a_n = (b_n)^n$, usar la \textbf{Prueba de la Raíz}.
    \item Si $a_n = f(n)$, donde $\int_{1}^{\infty} f(x)$ es fácilmente evaluable, usar la \textbf{Prueba de la Integral}.
\end{enumerate}

\section{Series de funciones}
Una serie de funciones es una suma infinita de funciones de la forma:
$$ \sum_{n=1}^{\infty} f_n(x) = f_1(x) + f_2(x) + f_3(x) + \dots $$
donde cada $f_n(x)$ es una función definida en un conjunto $A \subseteq \mathbb{R}$. \newline

Para cada valor de $x$ en $\sum_{n=1}^{\infty} f_n(x)$, se obtiene una serie numérica que puede diverger o converger.

\subsection{Suma parcial de una serie de funciones}
\begin{itemize}
    \item La $n$-ésima suma parcial de la serie de funciones $\sum_{n=1}^{\infty} f_n(x)$ está dada por:
    $$ S_n(x) = \sum_{i=1}^{n} f_i(x) = f_1(x) + f_2(x) + \dots + f_n(x) $$
    \item La suma de los infinitos términos de la serie es $S(x)$:
    $$S(x) = S_n(x) + R_n(x)$$
    donde $R_n(x)$ es el \textbf{residuo} o \textbf{error} al aproximar la suma de la serie por la suma parcial $S_n(x)$.
    \item El residuo $R_n(x)$ está dado por:
    $$ R_n(x) = \sum_{i=n+1}^{\infty} f_i(x) $$
    \item La sucesión de sumas parciales $\{S_n(x)\}$ define una función $S(x)$ llamada \textbf{suma de la serie de funciones} si el límite:
    $$ S(x) = \lim_{n \to \infty} S_n(x) $$
    existe para cada $x$ en el dominio.
\end{itemize}

\begin{tcolorbox}
    Esto implica que:
$$\lim_{n \to \infty} R_n(x) = 0$$
para todos los valores de $x$ para los que la serie \textbf{converge}. \newline

El conjunto de valores numéricos de $x$ para los cuales la serie de funciones converge se llama \textbf{dominio de convergencia}.
\end{tcolorbox}

\subsection{Convergencia de series de funciones}
\begin{itemize}
    \item \textbf{Convergencia puntual:} La serie $\sum_{n=1}^{\infty} f_n(x)$ converge puntualmente en un conjunto $A$ si para cada $x$ en $A$, la serie numérica $\sum_{n=1}^{\infty} f_n(x)$ converge.
    \item \textbf{Convergencia uniforme:} Si se considera una serie numérica convergente de términos positivos $\sum_{n=1}^{\infty} a_n$ tal que:
    $$ |f_n(x)| \leqslant a_n, \quad \forall n \in \mathbb{N}, \forall x \in A $$
    la serie $\sum_{n=1}^{\infty} f_n(x)$ converge uniformemente en cierto dominio $[x_1; x_2]$.
\end{itemize}

\subsection{Continuidad en series de funciones}
Si $\sum_{n=1}^{\infty} f_n(x)$ es una serie con funciones continuas de $x$ y la serie converge en un intervalo $[a; b]$, entonces la suma de un número finito de funciones continuas es otra función continua. \newline
No siempre la suma de un número infinito de funciones continuas es continua.

\subsection{Derivación e integración término a término}
Si la serie de potencias $\sum c_n (x - a)^n$ tiene radio de convergencia $R > 0$, entonces la función $f$ definida por
$$ f(x) = \sum_{n=0}^{\infty} c_n (x - a)^n $$
es \textbf{derivable} (y, por lo tanto, continua) en el intervalo $(a-R, a+R)$, y
\begin{itemize}
    \item $$f'(x) = \sum_{n=1}^{\infty} n c_n (x - a)^{n-1} $$
    \item $$\int f(x) dx = C + \sum_{n=0}^{\infty} c_n \frac{(x - a)^{n+1}}{n+1} $$
\end{itemize}

% Agregar a mano p6 y 7 PPT series parte b

\section{Series de Potencias}

\subsection{Series de potencias centradas en 0}
La forma $\sum_{n=0}^{\infty} a_nx^n$ es una serie de potencias centrada en 0. 

\subsubsection{Calculo del radio de convergencia - Cauchy-Hadamard}
Para toda serie de potencias existe un radio de convergencia $R$ de la serie de potencias $\sum_{n=0}^{\infty} a_nx^n$, siendo $0 \leqslant R \leqslant \infty$, tal que:
\begin{itemize}
    \item La serie converge absolutamente para $|x| < R$.
    \item La serie diverge para $|x| > R$.
    \item Para $|x| = R$, la serie puede converger o divergir. Cada extremo debe ser evaluado por separado.
\end{itemize}

El radio de convergencia $R$ deambas series de potencias está dado por:
$$ R = \frac{1}{\lim_{n \to \infty} \sqrt[n]{|a_n|}} $$
o, alternativamente, si el límite existe,
$$ R = \frac{1}{\lim_{n \to \infty} \sqrt[n]{|a_n|}} $$

\begin{tcolorbox}
    Si el límite del denominador es 0, entonces $R = \infty$ (la serie converge $\forall x$). Si el límite es infinito, entonces $R = 0$ (la serie solo converge en $x = 0$).
\end{tcolorbox}

\subsubsection{Propiedades de las series de potencias}
\begin{itemize}
    \item Si la serie de potencias $\sum_{n=0}^{\infty} a_n x^n$ es convergente para $x = x_1 (x_1 \neq 0)$, entonces es absolutamente convergente para todo $\frac{x}{|x|} < |x_1|$.
    \item Si la serie de potencias $\sum_{n=0}^{\infty} a_n x^n$ es divergente para $x = x_2 (x_2 \neq 0)$, entonces es divergente para todo $|x| > |x_2|$.
\end{itemize}

\begin{tcolorbox}
    La serie de potencias será \textbf{absolutamente convergente} en el intervalo $(-R, R)$ y \textbf{divergente} fuera de este intervalo.
\end{tcolorbox}

\subsection{Series de potencias centras en a}

Es una serie de la forma $\sum_{n=0}^{\infty} c_n (x - a)^n$, donde:
\begin{enumerate}
    \item La serie \textbf{converge} solo en $x = a$. $R = 0$ y el intervalo consta de un solo punto $\{a\}$.
    \item La serie \textbf{converge} para todos los valores de $x$. $R = \infty$ y el intervalo es $(-\infty, \infty)$.
    \item La serie \textbf{converge} para $|x - a| < R$ y \textbf{diverge} para $|x - a| > R$, donde $R$ es el \textbf{radio de convergencia}. La convergencia en los extremos del intervalo debe ser verificada por separado. Los intervaos pueden ser: 
    \begin{multicols}{2}
        \begin{itemize}
            \item $(a-R, a+R)$
            \item $(a-R, a+R]$
            \item $[a-R, a+R)$
            \item $[a-R, a+R]$
        \end{itemize}
    \end{multicols}
\end{enumerate}

\begin{tcolorbox}
    \textbf{Radio de convergencia ($R$):} Distancia desde el centro $a$ donde la serie converge.
    \textbf{Intervalo de convergencia:} Conjunto de valores de $x$ para los cuales la serie converge.
\end{tcolorbox}

\section{Series de Taylor y Maclaurin}
Si una función $f$ puede representarse mediante una serie de potencias centrada en $a$, es decir:
$$ f(x) = \sum_{n=0}^{\infty}c_n (x-a)^n $$
entonces sus coeficientes están determinados unívocamente por las derivadas sucesivas de la función evaluadas en el centro: $c_n = \frac{f^{(n)}(a)}{n!}$. La serie resultante:
$$ f(x) = \sum_{n=0}^{\infty} \frac{f^{(n)}(a)}{n!} (x-a)^n = f(a) + f'(a)(x-a) + \frac{f''(a)}{2!}(x-a)^2 + \dots $$
se llama \textbf{serie de Taylor} de $f$ centrada en $a$. Esta serie converge a $f(x)$ siempre que el resto tienda a cero dentro de su radio de convergencia $R$.

Para el caso especial en que el centro es el origen ($a = 0$), la serie se denomina \textbf{serie de Maclaurin}:
$$ f(x) = \sum_{n=0}^{\infty} \frac{f^{(n)}(0)}{n!} x^n = f(0) + f'(0)x + \frac{f''(0)}{2!}x^2 + \dots $$

\subsection{Resto de Taylor y Convergencia}
Se define el resto (o error de aproximación) como $R_n(x) = f(x) - T_n(x)$, donde $T_n(x)$ es el polinomio de Taylor de grado $n$. La igualdad $f(x) = \sum \frac{f^{(n)}(a)}{n!}(x-a)^n$ es válida si y solo si:
$$\lim_{n \to \infty} R_n(x) = 0 $$
para $|x-a| < R$.

\subsection{Desigualdad de Taylor}
Para estimar la precisión de la aproximación, si se puede garantizar que la derivada $(n+1)$-ésima está acotada, es decir, $|f^{(n+1)}(x)| \leqslant M$ para $x$ en un intervalo que contiene a $a$ y $d$, entonces el resto de Taylor satisface:
$$ |R_n(x)| \leqslant \frac{M}{(n+1)!} |x-a|^{n+1} $$

\begin{tcolorbox}[colback=blue!5,colframe=blue!75!black]
    \begin{itemize}
        \item $\lim_{n \to \infty} \frac{x^n}{n!} = 0$ para todo $x \in \mathbb{R}$.
        \item $e^x = \sum_{n=0}^{\infty} \frac{x^n}{n!}$ para todo $x$. En particular, si $x=1$, $e = \sum_{n=0}^{\infty} \frac{1}{n!}$.
    \end{itemize}
\end{tcolorbox}

\subsection{Serie binomial}
Si $k \in \mathbb{R}$ y $|x| < 1$, entonces
$$(1+x)^k = \sum_{n=0}^{\infty} \frac{k!}{n!(k-n)!} x^n$$

\subsection{Series importantes de Maclaurin y sus radios de convergencia}
\begin{table}[H]
    \centering
    \begin{tabular}{lll}
        \toprule
        Función & Serie de Maclaurin & Radio de Convergencia \\
        \midrule
        $e^x$ & $\sum_{n=0}^{\infty} \frac{x^n}{n!}$ & $\infty$ \\
        $\sin x$ & $\sum_{n=0}^{\infty} (-1)^n \frac{x^{2n+1}}{(2n+1)!}$ & $\infty$ \\
        $\cos x$ & $\sum_{n=0}^{\infty} (-1)^n \frac{x^{2n}}{(2n)!}$ & $\infty$ \\
        $\arctan x$ & $\sum_{n=0}^{\infty} (-1)^n \frac{x^{2n+1}}{2n+1}$ & $1$ \\
        $\frac{1}{1-x}$ & $\sum_{n=0}^{\infty} x^n$ & $1$ \\
        $\ln(1+x)$ & $\sum_{n=1}^{\infty} (-1)^{n-1} \frac{x^n}{n}$ & $1$ \\
        $(1+x)^k$ & $\sum_{n=0}^{\infty} \frac{k(k-1)\cdots(k-n+1)}{n!} x^n$ & $1$ \\
        \bottomrule
    \end{tabular}
\end{table}

%%---------------------------------------------------------------------------------------------


\end{document}